\documentclass[11pt]{article}
\usepackage[margin=1in]{geometry}
\usepackage{amsmath,amssymb,amsthm}
\usepackage{microtype}
\usepackage[hidelinks]{hyperref}
\newcommand{\R}{\mathbb R}
\newcommand{\Mat}{M_6(\R)}
\newcommand{\GL}{\mathrm{GL}_6(\R)}
\newcommand{\sgn}{\operatorname{sgn}}
\newcommand{\tr}{\operatorname{tr}}
\newcommand{\diag}{\operatorname{diag}}
\newtheorem{theorem}{Theorem}
\newtheorem{lemma}{Lemma}
\title{A piecewise quadratic function without a polynomial lattice representation}
\author{}
\date{}
\begin{document}
\maketitle

\begin{abstract}
We construct a continuous function on \(\R^{72}\) with a finite closed
semialgebraic cover on which its labels are homogeneous quadratic
polynomials, and prove that it has no finite lattice expression in
real polynomials of arbitrary degrees. Positive homogeneity imposes
a necessary finite condition on signs of linear and quadratic
polynomials. Pairs of matrices with opposite determinant signs
violate that condition. The construction concerns the unrestricted
Pierce--Birkhoff assertion; its displayed partition uses a determinant
of degree six.
\end{abstract}

A finite polynomial lattice expression is obtained from finitely many
real polynomials by finitely many applications of maximum and minimum.
By distributivity, such expressions are precisely finite maxima of
finite minima of polynomials. A continuous function is piecewise
polynomial here if a finite collection of closed semialgebraic sets
covers its whole domain and the function agrees with a polynomial
on each set.

Identify \(\Mat\times\Mat\) with \(\R^{72}\), using independent matrix
coordinates \(X,Y\). For \(1\le i\le6\) and signs
\(\sigma=(\sigma_j)_{j\ne i}\in\{-1,1\}^5\), put
\[
 q_{i,\sigma}(X,Y)=(XY)_{ii}-\sum_{j\ne i}\sigma_j(XY)_{ij},
 \qquad
 h(X,Y)=\min_{i,\sigma}q_{i,\sigma}(X,Y).
\]
There are \(192\) homogeneous quadratic forms \(q_{i,\sigma}\), all
with integer coefficients. Define
\begin{equation}\label{eq:function}
 f(X,Y)=
 \begin{cases}
 h(X,Y),&h(X,Y)>0\ \text{and}\ \det X>0,\\
 0,&\text{otherwise}.
 \end{cases}
\end{equation}

\begin{theorem}\label{thm:main}
The function \(f:\R^{72}\to\R\) in \eqref{eq:function} is continuous
and has a finite closed semialgebraic polynomial cover. It has no
finite lattice expression in real polynomials, with no restriction
on their degrees.
\end{theorem}

\section{The whole-space polynomial cover}

Writing \(P=XY\), minimization over the row signs gives
\[
 h(X,Y)=\min_i\left(P_{ii}-\sum_{j\ne i}|P_{ij}|\right).
\]
If \(h>0\), then \(P\) has positive diagonal and is strictly diagonally
dominant by rows. To see that it is invertible, suppose \(Pv=0\) with
\(v\ne0\), and choose \(i\) maximizing \(|v_i|>0\). The row equation gives
\[
 P_{ii}|v_i|
 \le \sum_{j\ne i}|P_{ij}|\,|v_j|
 \le \left(\sum_{j\ne i}|P_{ij}|\right)|v_i|
 <P_{ii}|v_i|,
\]
a contradiction. In particular, \(\det X\ne0\) on \(\{h>0\}\).

On that open set the sign of \(\det X\) is locally constant, so \(f\)
is continuous there. It is locally zero on \(\{h<0\}\). At \(h=0\),
its value is zero and
\[
 0\le f\le\max(h,0),
\]
which proves continuity at all remaining points, including singular
\(X\) and \(Y\).

Enumerate the quadratics as \(q_1,\ldots,q_{192}\), and define
\[
 F_0=\{\det X\le0\}\ \cup\ \bigcup_j\{q_j\le0\},
\]
\[
 F_i=\{\det X\ge0,\ q_j\ge0\text{ for all }j,\
                    q_i\le q_j\text{ for all }j\},
 \qquad 1\le i\le192.
\]
These \(193\) closed semialgebraic sets cover the whole space. Use
label \(0\) on \(F_0\) and label \(q_i\) on \(F_i\). On \(F_i\),
the minimum \(h=q_i\) is nonnegative. If it is positive, invertibility
and \(\det X\ge0\) give \(\det X>0\), so the label equals \(f\).
On \(F_i\cap F_0\), either a \(q_j\) is zero or \(\det X=0\);
in the latter case strict positivity of all \(q_j\) is impossible.
Thus \(q_i=0\) on this overlap. On \(F_i\cap F_j\), both labels
equal the common minimum. This verifies the asserted polynomial cover.

For every \(t>0\), simultaneous scaling gives
\begin{equation}\label{eq:homogeneous}
 f(tX,tY)=t^2f(X,Y).
\end{equation}
Indeed, each \(q_{i,\sigma}\) scales by \(t^2\), whereas
\(\det(tX)=t^6\det X\).

\section{A necessary condition from positive homogeneity}

\begin{lemma}\label{lem:sign}
Let \(g:\R^N\to\R\) satisfy \(g(tz)=t^2g(z)\) for every \(t>0\).
If \(g\) has a finite lattice expression in arbitrary real polynomials,
there is a finite family of polynomials of degree at most two such
that any two points sharing all signs in that family have the same
sign of \(g\).
\end{lemma}

\begin{proof}
Write \(g=\mathcal L(p_1,\ldots,p_m)\), where \(\mathcal L\) is a
fixed finite expression using only maximum and minimum. Decompose
each leaf into its homogeneous parts:
\[
 p_j=c_j+\ell_j+Q_j+\sum_{d\ge3}P_{j,d}.
\]
Here \(\ell_j\) has degree one and \(Q_j\) degree two; either can
be zero. Let \(\mathcal T=\{\ell_j,Q_j:1\le j\le m\}\).

For a fixed point \(z\), positive scaling commutes with the order
operations, and hence for every \(t>0\),
\[
 g(z)=\mathcal L\left(\frac{p_1(tz)}{t^2},\ldots,
                    \frac{p_m(tz)}{t^2}\right).
\]
Each entry has an extended-real limit \(u_j(z)\) as \(t\downarrow0\):
\[
 u_j(z)=
 \begin{cases}
 \sgn(c_j)\,\infty,&c_j\ne0,\\
 \sgn(\ell_j(z))\,\infty,&c_j=0,\ \ell_j(z)\ne0,\\
 Q_j(z),&c_j=0,\ \ell_j(z)=0.
 \end{cases}
\]
Terms of degree at least three tend to zero after division by \(t^2\)
at each fixed \(z\). Finite minimum and maximum are continuous on
the extended real line. For example, conjugation by the increasing
homeomorphism \((2/\pi)\arctan:[-\infty,\infty]\to[-1,1]\)
reduces this to their ordinary continuity. No addition of infinities
occurs. Therefore
\begin{equation}\label{eq:limit}
 g(z)=\mathcal L(u_1(z),\ldots,u_m(z)).
\end{equation}

Extend sign by \(\sgn(-\infty)=-1\) and \(\sgn(+\infty)=1\).
Every nondecreasing map between totally ordered sets preserves
minimum and maximum. Applying sign to the exact equality
\eqref{eq:limit} gives
\[
 \sgn g(z)=\mathcal L(\sgn u_1(z),\ldots,\sgn u_m(z)).
\]
The inputs on the right are determined by the fixed constants \(c_j\)
and the signs of \(\ell_j(z),Q_j(z)\). Thus equal joint signs for
\(\mathcal T\) imply equal signs of \(g\).

The limit is taken before sign is applied; sign is not interchanged
with a limit at zero. Only pointwise convergence is used.
\end{proof}

\section{All quadratic equations on an orbit}

Put
\[
 D=\diag(1,1,0,0,0,0),\qquad
 E=\diag(0,0,1,1,0,0),
\]
and consider
\[
 \mathcal O=\{(UDV^{-1},VEU^{-1}):U,V\in\GL\}.
\]
Its pairs have ranks two and two, and satisfy \(XY=YX=0\).

\begin{lemma}\label{lem:orbit}
A real polynomial of total degree at most two vanishes on
\(\mathcal O\) if and only if it is a constant real linear
combination of the entries of \(XY\) and \(YX\).
\end{lemma}

\begin{proof}
The reverse implication is immediate. For nonzero real \(a,b\),
replace \(U\) by
\[
 U\diag(a,a,b^{-1},b^{-1},1,1)
\]
and keep \(V\) fixed. This replaces \((X,Y)\) by \((aX,bY)\).
Consequently, separating coefficients in \(a,b\), each bihomogeneous
component of a polynomial vanishing on \(\mathcal O\) vanishes
separately.

The \(X\)-projection and the \(Y\)-projection of \(\mathcal O\)
are both the entire rank-two matrix locus, by independent changes
of basis on the left and right. Its closure contains every rank-one
matrix and zero. A homogeneous linear polynomial vanishing there
is zero by evaluation at matrix units. For a homogeneous quadratic,
these evaluations eliminate the squared-coordinate coefficients,
and evaluation at sums of two distinct matrix units eliminates all
mixed coefficients. Each such sum has rank at most two. The
constant component is also zero. Thus only a bilinear polynomial
\(B(X,Y)\) can remain.

We next show that every nonzero rank-one pair
\begin{equation}\label{eq:rankone}
 X=uv^T,\qquad Y=st^T,\qquad t^Tu=0,\quad v^Ts=0
\end{equation}
lies in the Euclidean closure of \(\mathcal O\).
Regard \(X\) as a map from a second vector space to a first, and
\(Y\) in the reverse direction. In the first space choose a basis
with first vector \(u\), a second vector on which \(t^T\) is one,
and all other vectors in \(\ker(t^T)\). In the second space choose
a first vector on which \(v^T\) is one, second vector \(s\), and
all other vectors in \(\ker(v^T)\). The two zero-pairing conditions
make these choices possible. In these bases \(X=E_{11}\) and
\(Y=E_{22}\), where \(E_{ij}\) are matrix units.

Let \(P\) interchange coordinates two and three. For \(\delta\ne0\),
set
\[
 U_\delta=P\diag(1,\delta,1,\delta^{-1},1,1),\qquad V_\delta=P.
\]
Directly,
\[
 U_\delta D V_\delta^{-1}=E_{11}+\delta E_{33},\qquad
 V_\delta E U_\delta^{-1}=E_{22}+\delta E_{44}.
\]
These tend to the normalized rank-one pair. Composing with the
two fixed changes of basis proves the closure assertion. Hence
\(B\) vanishes on all pairs \eqref{eq:rankone}; if one matrix is
zero it vanishes by bilinearity.

Write
\[
 B(X,Y)=\sum_{i,j,k,l}b_{ij,kl}X_{ij}Y_{kl},
 \qquad
 T(A,C)=\sum_{i,j,k,l}b_{ij,kl}A_{il}C_{kj}.
\]
The second expression defines a bilinear form on two matrix spaces,
and
\begin{equation}\label{eq:tensor}
 B(uv^T,st^T)=T(ut^T,sv^T).
\end{equation}
The two arguments on the right range independently over all rank-one
traceless matrices: their trace conditions are exactly
\(t^Tu=0\) and \(v^Ts=0\). Such matrices span the entire traceless
matrix space. They include all off-diagonal units, and
\[
 E_{ii}-E_{jj}
 =(e_i+e_j)(e_i-e_j)^T+E_{ij}-E_{ji}
\]
supplies all diagonal differences. Bilinearity therefore shows that
\(T\) vanishes on the product of the two full traceless spaces.

Decomposing each argument into a traceless matrix plus its trace
times \(I/6\) now gives
\[
 T(A,C)=\tr(A)L(C)+\tr(C)M(A),
\]
where one can take
\[
 L(C)=\frac{T(I,C)}6-\frac{\tr(C)T(I,I)}{36},
 \qquad M(A)=\frac{T(A,I)}6.
\]
Substituting in \eqref{eq:tensor}, the first summand is \(L(YX)\),
and the second is \(M(XY)\). The trace formula holds for all \(A,C\),
so this gives
\(B(X,Y)=L(YX)+M(XY)\) for arbitrary rank-one \(X,Y\), without
requiring their products to vanish. Both sides are bilinear and
rank-one matrices span the full matrix space. The identity therefore
holds for all \(X,Y\), as required.
\end{proof}

\section{Pairs with identical signs for any finite quadratic family}

\begin{lemma}\label{lem:pairs}
For any finite family of real polynomials \(p_1,\ldots,p_M\)
of degree at most two, there are two actual real matrix pairs
\(z_+,z_-\) with identical signs for every \(p_j\), but
\(f(z_+)>0\) and \(f(z_-)=0\).
\end{lemma}

\begin{proof}
For each \(p_j\) not identically zero on \(\mathcal O\), the function
\[
 (U,V)\longmapsto p_j(UDV^{-1},VEU^{-1})
\]
is a nonzero rational function of the entries of \(U,V\).
Multiplication by \((\det U\det V)^2\) clears its denominators
and gives a nonzero polynomial numerator. The product of these
finitely many nonzero numerators is nonzero; an empty product is one.

A nonzero real polynomial cannot vanish on a nonempty Euclidean
open set. Indeed, induction on the number of variables reduces
vanishing on an open box to the univariate root bound by comparing
coefficient polynomials. Choose an open neighborhood of \((I,I)\)
where both determinants are positive. There are \(U,V\) in this
neighborhood at which all the indicated numerators are nonzero.
Thus every \(p_j\) not vanishing on \(\mathcal O\) is nonzero at
the common point
\[
 o=(UDV^{-1},VEU^{-1}).
\]

Fix these \(U,V\). For \(\epsilon>0\) and \(\eta\in\{-1,1\}\), put
\[
 D_X^\eta=\diag(1,1,\epsilon^2,\epsilon^2,\epsilon,\eta\epsilon),
 \qquad
 D_Y^\eta=\diag(\epsilon^2,\epsilon^2,1,1,\epsilon,\eta\epsilon),
\]
\[
 X_\eta=UD_X^\eta V^{-1},\qquad
 Y_\eta=VD_Y^\eta U^{-1}.
\]
Both pairs tend to \(o\) as \(\epsilon\downarrow0\). Exact
multiplication and determinants give
\begin{equation}\label{eq:products}
 X_\eta Y_\eta=Y_\eta X_\eta=\epsilon^2I,
 \qquad
 \det X_\eta=\frac{\det U}{\det V}\,\eta\epsilon^6.
\end{equation}
If \(p_j\) vanishes on \(\mathcal O\), Lemma~\ref{lem:orbit}
and \eqref{eq:products} give exactly equal values on both pairs.
If \(p_j(o)\ne0\), continuity gives the same strict sign on both
pairs for all sufficiently small positive \(\epsilon\).
There are finitely many such requirements, so a single positive
real \(\epsilon\) makes every sign agree.

All \(192\) forms \(q_{i,\sigma}\) equal \(\epsilon^2\) on both
pairs. Since \(\det U/\det V>0\), their \(X\) determinants have
opposite signs. Therefore
\[
 f(X_+,Y_+)=\epsilon^2>0,\qquad f(X_-,Y_-)=0.
\]
The choices are made in the order: the finite polynomial family,
then \(U,V\), then \(\epsilon\). The function \(f\) is fixed throughout.
\end{proof}

\section{Conclusion}

\begin{proof}[Proof of Theorem~\ref{thm:main}]
Section~1 verifies continuity and the finite closed polynomial
cover on all of \(\R^{72}\).
If \(f\) had any finite polynomial lattice expression,
positive homogeneity \eqref{eq:homogeneous} and Lemma~\ref{lem:sign}
would make its sign determined by some finite family of
degree-at-most-two polynomials. Lemma~\ref{lem:pairs} supplies
two real points with identical signs for that family and different
signs of \(f\), a contradiction.
\end{proof}

The displayed semialgebraic cover uses a determinant of degree six.
The theorem does not assert that \(f\) has one shared polynomial
label on each entire sign realization of some finite family of
quadratic partition polynomials. That additional partition requirement
is different from the unrestricted whole-space assertion disproved
here.

\end{document}

